% GENERATED FILE. Edit canonical lessons, then run npm run generate:books.
% canonical-source-hash: fnv1a64:6248b02e4d4f8ef8
% canonical-lessons: MR-C514-karya, MR-C514-kamkaj, MR-C514-pali, MR-C514-urakne, MR-C514-dhadpadne

\chapter{Getting It Done}
\label{ch:mr-getting-it-done}
\hlchaptermodality{514}

\begin{chapteropening}
\textbf{By the end of this chapter:} \emph{I can talk about tasks, shifts and finishing work.}

Complete the last lesson of chapter 514: I can talk about tasks, shifts and finishing work.
\end{chapteropening}

\section[\texorpdfstring{\mr{कार्य} (kārya)}{कार्य (kārya)}]{\mr{कार्य} (kārya) — a task; a function}

\label{lesson:MR-C514-karya}



\begin{quote}
Before the new one: say the Marathi for to present, then the Marathi for to host a programme.
\end{quote}

\subsection*{You'll want to know: \mr{कार्य}}
\textbf{\mr{कार्य}} — \emph{kārya} — \textquotedblleft{}a task; a function\textquotedblright{}.

\begin{quote}
\mr{हे} \mr{कार्य} \mr{आज} \mr{संध्याकाळपर्यंत} \mr{पूर्ण} \mr{झालं}.
\end{quote}

\begin{quote}
\emph{he kārya āj sandhyākāḷparyant pūrṇa jhālaṁ.}
\end{quote}

\begin{quote}
This task was finished by this evening.
\end{quote}

Neuter.

\begin{cousinweb}
\textbf{\mr{कार्य}} (\emph{kārya}) is a tatsama, taken from Sanskrit \emph{kārya}, \textquotedblleft{}what is to be done, work\textquotedblright{} (Monier-Williams).
\end{cousinweb}

\begin{grammarlens}[title={What you've built}]
One more word for describing what you lived and felt.
\end{grammarlens}

\subsection*{Your turn}
\begin{itemize}
\raggedright
  \item \emph{Say it:} \emph{kārya}
  \item \emph{Say it:} \emph{kārya}, once more
  \item \emph{Say it:} \emph{he kārya āj sandhyākāḷparyant pūrṇa jhālaṁ.} — This task was finished by this evening.
  \item \emph{Recall:} say the Marathi for to present, then the Marathi for to host a programme, then say \emph{kārya} again
\end{itemize}

\begin{tcolorbox}[breakable,colback=teal!4,colframe=teal!35!black,title={Before you move on}]
What does \mr{कार्य} mean? (\textquotedblleft{}A task; a function\textquotedblright{}.) Say it once more.
\end{tcolorbox}
\section[\texorpdfstring{\mr{कामकाज} (kāmkāj)}{कामकाज (kāmkāj)}]{\mr{कामकाज} (kāmkāj) — the day's business; proceedings}

\label{lesson:MR-C514-kamkaj}



\begin{quote}
Before the new one: say the Marathi for to host a programme, then the Marathi for a task.
\end{quote}

\subsection*{You'll want to know: \mr{कामकाज}}
\textbf{\mr{कामकाज}} — \emph{kāmkāj} — \textquotedblleft{}the day's business; proceedings\textquotedblright{}.

\begin{quote}
\mr{सुट्टीमुळे} \mr{कार्यालयाचं} \mr{कामकाज} \mr{बंद} \mr{होतं}.
\end{quote}

\begin{quote}
\emph{suṭṭīmuḷe kāryālayācaṁ kāmkāj band hotaṁ.}
\end{quote}

\begin{quote}
Because of the holiday the office's business was closed.
\end{quote}

Neuter.

\begin{cousinweb}
\textbf{\mr{कामकाज}} (\emph{kāmkāj}) is a Marathi pairing of \mr{काम}, \textquotedblleft{}work\textquotedblright{}, and \mr{काज}, \textquotedblleft{}business\textquotedblright{}, which is inherited from Sanskrit \emph{kārya} (CDIAL 3078).
\end{cousinweb}

\begin{grammarlens}[title={What you've built}]
One more word for describing what you lived and felt.
\end{grammarlens}

\subsection*{Your turn}
\begin{itemize}
\raggedright
  \item \emph{Say it:} \emph{kāmkāj}
  \item \emph{Say it:} \emph{kāmkāj}, once more
  \item \emph{Say it:} \emph{suṭṭīmuḷe kāryālayācaṁ kāmkāj band hotaṁ.} — Because of the holiday the office's business was closed.
  \item \emph{Recall:} say the Marathi for to host a programme, then the Marathi for a task, then say \emph{kāmkāj} again
\end{itemize}

\begin{tcolorbox}[breakable,colback=teal!4,colframe=teal!35!black,title={Before you move on}]
What does \mr{कामकाज} mean? (\textquotedblleft{}The day's business; proceedings\textquotedblright{}.) Say it once more.
\end{tcolorbox}
\section[\texorpdfstring{\mr{पाळी} (pāḷī)}{पाळी (pāḷī)}]{\mr{पाळी} (pāḷī) — a shift; a turn}

\label{lesson:MR-C514-pali}



\begin{quote}
Before the new one: say the Marathi for a task, then the Marathi for the day's business.
\end{quote}

\subsection*{You'll want to know: \mr{पाळी}}
\textbf{\mr{पाळी}} — \emph{pāḷī} — \textquotedblleft{}a shift; a turn\textquotedblright{}.

\begin{quote}
\mr{माझ्या} \mr{भावाची} \mr{रात्रीची} \mr{पाळी} \mr{असते}.
\end{quote}

\begin{quote}
\emph{mājhyā bhāvācī rātrīcī pāḷī aste.}
\end{quote}

\begin{quote}
My brother works the night shift.
\end{quote}

Feminine.

\begin{cousinweb}
\textbf{\mr{पाळी}} (\emph{pāḷī}) is inherited: CDIAL lists Marathi \emph{pāḷī}, \textquotedblleft{}a turn\textquotedblright{}, under \emph{*pāḍi}, \textquotedblleft{}row, line, turn\textquotedblright{} (CDIAL 8041).
\end{cousinweb}

\begin{grammarlens}[title={What you've built}]
One more word for describing what you lived and felt.
\end{grammarlens}

\subsection*{Your turn}
\begin{itemize}
\raggedright
  \item \emph{Say it:} \emph{pāḷī}
  \item \emph{Say it:} \emph{pāḷī}, once more
  \item \emph{Say it:} \emph{mājhyā bhāvācī rātrīcī pāḷī aste.} — My brother works the night shift.
  \item \emph{Recall:} say the Marathi for a task, then the Marathi for the day's business, then say \emph{pāḷī} again
\end{itemize}

\begin{tcolorbox}[breakable,colback=teal!4,colframe=teal!35!black,title={Before you move on}]
What does \mr{पाळी} mean? (\textquotedblleft{}A shift; a turn\textquotedblright{}.) Say it once more.
\end{tcolorbox}
\section[\texorpdfstring{\mr{उरकणे} (urakṇe)}{उरकणे (urakṇe)}]{\mr{उरकणे} (urakṇe) — to get done; to finish off}

\label{lesson:MR-C514-urakne}



\begin{quote}
Before the new one: say the Marathi for the day's business, then the Marathi for a shift.
\end{quote}

\subsection*{You'll want to know: \mr{उरकणे}}
\textbf{\mr{उरकणे}} — \emph{urakṇe} — \textquotedblleft{}to get done; to finish off\textquotedblright{}.

\begin{quote}
\mr{दुपारच्या} \mr{आधी} \mr{त्याने} \mr{सगळी} \mr{कामं} \mr{उरकली}.
\end{quote}

\begin{quote}
\emph{dupārcyā ādhī tyāne sagḷī kāmaṁ uraklī.}
\end{quote}

\begin{quote}
He got all the jobs done before noon.
\end{quote}

Past, with -\mr{ने} (he …): \mr{त्याने} \mr{काम} \mr{उरकलं}, \emph{tyāne kām uraklaṁ}.

With a plural neuter object the past ends -\mr{ली}: \mr{कामं} \mr{उरकली}.

\begin{cousinweb}
\textbf{\mr{उरकणे}} (\emph{urakṇe}) is inherited: CDIAL lists Marathi \emph{urakṇẽ}, \textquotedblleft{}to be over, be finished\textquotedblright{}, under \emph{*urvara}, \textquotedblleft{}left over\textquotedblright{} (CDIAL 2356).
\end{cousinweb}

\begin{grammarlens}[title={What you've built}]
One more word for describing what you lived and felt.
\end{grammarlens}

\subsection*{Your turn}
\begin{itemize}
\raggedright
  \item \emph{Say it:} \emph{urakṇe}
  \item \emph{Say it:} \emph{urakṇe}, once more
  \item \emph{Say it:} \emph{dupārcyā ādhī tyāne sagḷī kāmaṁ uraklī.} — He got all the jobs done before noon.
  \item \emph{Recall:} say the Marathi for the day's business, then the Marathi for a shift, then say \emph{urakṇe} again
\end{itemize}

\begin{tcolorbox}[breakable,colback=teal!4,colframe=teal!35!black,title={Before you move on}]
What does \mr{उरकणे} mean? (\textquotedblleft{}To get done; to finish off\textquotedblright{}.) Say it once more.
\end{tcolorbox}
\section[\texorpdfstring{\mr{धडपडणे} (dhaḍpaḍṇe)}{धडपडणे (dhaḍpaḍṇe)}]{\mr{धडपडणे} (dhaḍpaḍṇe) — to struggle; to scramble}

\label{lesson:MR-C514-dhadpadne}



\begin{quote}
Before the new one: say the Marathi for a shift, then the Marathi for to get done.
\end{quote}

\subsection*{You'll want to know: \mr{धडपडणे}}
\textbf{\mr{धडपडणे}} — \emph{dhaḍpaḍṇe} — \textquotedblleft{}to struggle; to scramble\textquotedblright{}.

\begin{quote}
\mr{नवीन} \mr{कामात} \mr{सुरुवातीला} \mr{मी} \mr{खूप} \mr{धडपडलो}.
\end{quote}

\begin{quote}
\emph{navīn kāmāt suruvātīlā mī khūp dhaḍpaḍlo.}
\end{quote}

\begin{quote}
At first I struggled a lot in the new job.
\end{quote}

Past (he …): \mr{तो} \mr{धडपडला}, \emph{to dhaḍpaḍlā}.

A woman says \mr{धडपडले}.

\begin{cousinweb}
\textbf{\mr{धडपडणे}} (\emph{dhaḍpaḍṇe}) is an expressive Marathi verb built on the sound word \mr{धडपड}; its origin is not settled.
\end{cousinweb}

\begin{grammarlens}[title={What you've built}]
One more word for describing what you lived and felt.
\end{grammarlens}

\subsection*{Your turn}
\begin{itemize}
\raggedright
  \item \emph{Say it:} \emph{dhaḍpaḍṇe}
  \item \emph{Say it:} \emph{dhaḍpaḍṇe}, once more
  \item \emph{Say it:} \emph{navīn kāmāt suruvātīlā mī khūp dhaḍpaḍlo.} — At first I struggled a lot in the new job.
  \item \emph{Recall:} say the Marathi for a shift, then the Marathi for to get done, then say \emph{dhaḍpaḍṇe} again
\end{itemize}

\begin{tcolorbox}[breakable,colback=teal!4,colframe=teal!35!black,title={Before you move on}]
What does \mr{धडपडणे} mean? (\textquotedblleft{}To struggle; to scramble\textquotedblright{}.) Say it once more.
\end{tcolorbox}
